\documentclass[10pt,a4paper]{amsart}
\usepackage[margin=19mm]{geometry}
\usepackage{amsmath,amssymb,mathtools}
\usepackage[scr=rsfs,cal=euler]{mathalfa}
\usepackage{xfrac}
\usepackage[unicode,hidelinks,bookmarks=false]{hyperref}
\newcommand{\ie}{i.e.\ }
\newcommand{\cf}{cf.\ }
\newcommand{\resp}{respectively\ }
\newcommand{\Subsection}{\subsection}
\providecommand{\nobreakdash}{\nobreak\hbox{-}}
\setlength{\emergencystretch}{2em}
\multlinegap0pt
\usepackage{amssymb}
\usepackage{algorithm}\usepackage{algpseudocode}
\usepackage{enumerate}
\newtheorem{proposition}{Proposition}
\title{Update-Magnitude State Redistribution (UM-SRD):\\ A Shut-off Extension of Weighted SRD for Cut-Cell Methods}
\author{Justo E. Karell}
\date{Source: arXiv:2605.04863v2 (CC BY 4.0)}
\begin{document}
\maketitle

\noindent\emph{Source and licence.} Update-Magnitude State Redistribution (UM-SRD):\\ A Shut-off Extension of Weighted SRD for Cut-Cell Methods. Version arXiv:2605.04863v2 (CC BY 4.0), distributed by arXiv under the Creative Commons Attribution 4.0 International licence, http://creativecommons.org/licenses/by/4.0/, as recorded in the arXiv licence metadata for this exact version and shown on the abstract page https://arxiv.org/abs/2605.04863v2. Full licence notice: \texttt{licenses/arxiv-2605.04863-CC-BY-4.0.txt}.\par\smallskip

\noindent\textit{Editorial scope.} Counted unit: the article's proposition prop:local-consistency (source0.tex lines 443--450). The article's complete original proof of it is delivered with it (source0.tex lines 452--468, one \texttt{proof} environment of 17 lines). No mathematics is written, repaired or re-derived: the statement and the proof are the article's own lines, in the article's own notation, and no retained line is substituted or re-flowed.  Retained uncounted as the source-local context the proof needs: lines 259--262; lines 432--436.  Nothing was omitted from any retained span, so the package carries no omission record.  No retained line contains a citation, so the wrapper carries no bibliography and no bibliography entry.  No retained line contains a figure environment, an image-inclusion command or drawing code, so no image is delivered and the package declares no asset dependency.  Editorial formatting only, in addition: an AMS \texttt{amsart} preamble that restates the article's own declarations and macros used by the retained lines, namely the environments \texttt{proposition} on separate counters, together with the standard packages those lines need.  The retained environments print in this document as Proposition 1 in order of appearance, whereas the article prints the same items with the numbering of its own full document.  The complete native source of the article as distributed in the arXiv e-print for this exact version is bundled unchanged as \texttt{sources/199904/source0.tex}.
\begin{equation}
  \eta_j = \frac{\Delta U_{\max,j}}{\varepsilon + \Delta U_{\max,j}},
  \label{eq:eta}
\end{equation}

\begin{equation}
  U^{n+1,\mathrm{UM}} - U^{n+1,\mathrm{base}}
  = s_j^n (S_j - \mathrm{Id})\, U^{n+1,\mathrm{base}},
  \label{eq:perturbation}
\end{equation}

\begin{proposition}[First-order accuracy]
\label{prop:local-consistency}
With the normalized indicator of \eqref{eq:eta}, UM-SRD remains first-order 
accurate: for smooth solutions,
\[
  \|\tau^{n,\mathrm{UM}} - \tau^{n,\mathrm{base}}\|_\infty = \mathcal{O}(h).
\]
\end{proposition}

\begin{proof}
From \eqref{eq:perturbation}, $\tau^{n,\mathrm{UM}} - \tau^{n,\mathrm{base}} 
= s_j^n (S_j - \mathrm{Id}) B(u^n)$. In the one-small-cell merging-left 
configuration,
\begin{align*}
  [(S_j - \mathrm{Id})V]_{j-1} &= \alpha(V_j - V_{j-1}), \\
  [(S_j - \mathrm{Id})V]_j     &= (1-\alpha)(V_{j-1} - V_j),
\end{align*}
with $(S_j - \mathrm{Id})V = 0$ at all other indices. For smooth solutions 
$\|(S_j - \mathrm{Id})B(u^n)\|_\infty = \mathcal{O}(h)$, and since 
$s_j^n = \mathcal{O}(1)$,
\[
  \|\tau^{n,\mathrm{UM}} - \tau^{n,\mathrm{base}}\|_\infty
  = \mathcal{O}(1) \cdot \mathcal{O}(h) = \mathcal{O}(h), \qedhere
\]
matching the truncation error of the underlying first-order scheme.
\end{proof}

\end{document}
